% 08 - CAS LIMITES
% ==============================================================
\sectionintro{08 · Retour d'essais}{Cas limites détectés et corrigés}{Le relevé d'essais décrit trois comportements limites ; leurs corrections sont présentes dans le bloc fonctionnel actuel.}

\begin{card}[colback=ekwarning,colframe=orange!35]
  \textbf{1 · Consigne de pièces égale à zéro}\par
  \begin{tabularx}{\linewidth}{@{}L{27mm}Y@{}}
    Symptôme & \texttt{0 >= 0} rendait chaque bac plein immédiatement et lançait une rotation continue.\\
    Correction & contrôle de \texttt{nPieceBac <= 0} avant l'exécution du cycle \texttt{AUTO}.\\
    Résultat & passage en \texttt{DEFAUT} et neutralisation des deux commandes du vérin.
  \end{tabularx}
\end{card}

\begin{card}[colback=ekwarning,colframe=orange!35]
  \textbf{2 · Huit bacs déjà pleins}\par
  \begin{tabularx}{\linewidth}{@{}L{27mm}Y@{}}
    Symptôme & après le retour 8 vers 1, le programme pouvait continuer à parcourir les bacs sans intervention opérateur.\\
    Correction & ajout de l'état \texttt{ATTENTE\_BAC} lorsque le compteur du nouveau bac a déjà atteint la consigne.\\
    Résultat & arrêt sur le bac concerné, réinitialisation par l'opérateur, puis reprise du comptage.
  \end{tabularx}
\end{card}

\begin{card}[colback=ekwarning,colframe=orange!35]
  \textbf{3 · Temps vérin égal à zéro}\par
  \begin{tabularx}{\linewidth}{@{}L{27mm}Y@{}}
    Symptôme & le \texttt{TON} terminait immédiatement et les états pouvaient s'enchaîner à la cadence du cycle PLC.\\
    Correction & contrôle de \texttt{tVerin <= T\#0MS} avec passage en \texttt{DEFAUT}.\\
    Résultat & aucune commande de mouvement tant que le paramètre n'est pas corrigé et acquitté.
  \end{tabularx}
\end{card}

\begin{primarycard}
  \textbf{Apport principal des essais}\par\small
  Le relevé couvre le cycle nominal, les valeurs limites et la saturation des huit bacs. La revue actuelle confirme la présence des trois corrections dans les sources ; elle ne constitue pas un nouveau rejeu du runtime.
\end{primarycard}

\newpage

% ==============================================================
